\documentclass[11pt]{article}
\usepackage[utf8]{inputenc}
\usepackage[T1]{fontenc}
\usepackage{amsmath,amssymb}
\usepackage{booktabs}
\usepackage{graphicx}
\usepackage{hyperref}
\usepackage{geometry}
\usepackage{multirow}
\usepackage{xcolor}
\geometry{margin=1in}

\title{Lexical Confounds in Multi-Turn Jailbreak Detection:\\A Register-Normalization Audit}

\author{Anonymous}
\date{}

\begin{document}
\maketitle

\begin{abstract}
Multi-turn jailbreak detection benchmarks report high accuracy, but we show that these results are systematically inflated by lexical confounds. We audit a representative persuasion-grounded detection pipeline and find that a TF-IDF classifier achieves perfect separation (F1=1.0) on the raw data, indicating that attack and benign conversations occupy entirely disjoint lexical registers. We propose a two-round register-normalization protocol that reduces this confound from TF-IDF separation of 1.0 to 0.73. On normalized data, a persuasion-grounded detector (DeBERTa fine-tuned on 8-class persuasion taxonomy + GRU) achieves F1=0.86, only 1 percentage point above a vanilla semantic baseline (F1=0.85, $p=0.62$). Scrambled-label ablation confirms that persuasion strategy labels contribute no significant signal beyond regularization. Adversarial paraphrasing degrades both systems equally, and cross-attack generalization to unseen attack types (FITD, Deceptive Delight) yields F1 below 0.58 for both approaches. These findings demonstrate that (1) lexical shortcuts constitute a structural confound in multi-turn jailbreak evaluation, not merely a data hygiene issue, and (2) persuasion-grounded features, under current taxonomy granularity and data scale, provide no measurable robustness or generalization advantage over generic semantic representations. We release our register-normalization protocol as a diagnostic tool for future multi-turn safety benchmarks.
\end{abstract}

\section{Introduction}

Multi-turn jailbreak attacks distribute adversarial intent across a sequence of individually benign conversational turns, defeating single-turn content filters. Recent detection methods report high accuracy by training classifiers on multi-turn conversation data. However, few studies audit whether reported performance reflects genuine detection of adversarial strategy or merely exploitation of superficial distributional differences between attack and benign conversations.

This distinction matters for deployment. If detectors rely on lexical register differences (attack conversations use more formal, template-derived language) rather than structural features of adversarial manipulation, they will fail against attackers who simply match the register of benign users. The gap between benchmark performance and real-world robustness remains unmeasured for most published multi-turn detectors.

We investigate this question through a systematic audit of a persuasion-grounded detection pipeline. The persuasion grounding hypothesis posits that multi-turn attacks follow a characteristic persuasion playbook (rapport building, authority framing, escalation) that constitutes a structural necessity for attack success, and that encoding this pragmatic structure should yield detectors more robust than those relying on lexical or semantic features alone. We test this hypothesis under controlled conditions.

Our audit reveals three findings. First, the raw multi-turn jailbreak data exhibits perfect lexical separation (TF-IDF F1=1.0), meaning any classifier trained on this data may learn lexical shortcuts rather than adversarial structure. This is not a quirk of one dataset: the register gap arises because attack conversations are generated by prompt templates with distinctive vocabulary, while benign conversations come from naturalistic user interactions. Second, after applying our two-round register-normalization protocol, the persuasion-grounded detector provides no statistically significant advantage over a vanilla semantic baseline (1-2 pp, $p>0.3$). Third, scrambled persuasion labels degrade performance by only 2 pp, confirming that the taxonomy signal is not driving detection.

Our contributions are:
\begin{itemize}
\item A register-normalization protocol for auditing lexical confounds in multi-turn detection benchmarks, demonstrated to reduce TF-IDF separation from 1.0 to 0.73 in two rounds.
\item Empirical evidence that persuasion-grounded features (8-class taxonomy, DeBERTa encoder) provide no significant advantage over vanilla semantic baselines on properly controlled data (N=72 test set, all comparisons $p>0.3$).
\item A systematic evaluation revealing that both approaches fail equivalently under adversarial perturbation, cross-attack generalization, and hard-negative scenarios, suggesting that current multi-turn detection capabilities are overestimated.
\end{itemize}

\section{Related Work}

\paragraph{Multi-turn jailbreak attacks.}
Crescendo distributes adversarial intent across multiple rapport-building turns before requesting harmful content. Foot-in-the-Door (FITD) uses incremental compliance requests, while Deceptive Delight interleaves benign and harmful content. Bullwinkel et al. show that per-turn representations in multi-turn attacks are deliberately kept in a ``benign'' region of embedding space, defeating representation-level single-turn defenses.

\paragraph{Multi-turn detection.}
DeepContext models generic semantic drift across turns using contextual embeddings. Persuasion-based approaches ground detection in discourse-level strategy classification, motivated by Zeng et al.'s Persuasive Adversarial Prompts (PAP) taxonomy and Hagendorff's finding that identifiable persuasion strategies appear in 97\% of high-ASR attacks.

\paragraph{Evaluation methodology for safety classifiers.}
Prior work on single-turn toxicity detection has identified lexical confounds and annotation artifacts. Gonen and Goldberg show that debiasing methods often preserve cluster structure detectable by simple classifiers. We extend this line of inquiry to multi-turn jailbreak detection, identifying register-level confounds specific to the multi-turn setting.

\section{Method}

\subsection{Detection Architecture}

We evaluate a two-stage persuasion-grounded detector:

\textbf{Stage 1: Persuasion Encoder.} DeBERTa-v3-base is fine-tuned on a multi-task objective: 8-class persuasion strategy classification (rapport, authority appeal, logical appeal, emotional appeal, urgency, reciprocity, social proof, reframing) and binary intent classification (jailbreak vs. benign turn-level intent). The [CLS] embeddings from this encoder serve as pragmatic-aware turn representations.

\textbf{Stage 2: Sequence Classifier.} A bidirectional GRU processes the sequence of turn embeddings and produces a conversation-level binary classification (jailbreak vs. benign).

\textbf{Baseline.} The identical GRU architecture processes vanilla DeBERTa-v3-base [CLS] embeddings (no persuasion fine-tuning).

\textbf{Additional baselines.} TF-IDF + Logistic Regression operates on concatenated turn text. A DeepContext proxy uses BERT-base + GRU with the same sequence architecture.

\subsection{Register-Normalization Protocol}

We propose a two-round protocol to diagnose and mitigate lexical confounds:

\textbf{Round 1: Template paraphrasing.} All conversations are paraphrased using an instruction-following LLM with the directive to preserve semantic content while removing distinctive register markers (formality level, template-specific phrases, stylistic flourishes). TF-IDF separation is measured after paraphrasing.

\textbf{Round 2: Targeted domain-term normalization.} Remaining high-TF-IDF features are identified (terms with class-conditional frequency ratio $>3$). These domain-specific terms are replaced with generic equivalents or distributed uniformly across both classes through targeted rewriting.

\textbf{Gate criterion.} Data passes the audit gate when TF-IDF + Logistic Regression achieves F1 $< 0.75$ on the normalized data, indicating that lexical features alone cannot near-perfectly separate classes.

\subsection{Data}

Training data consists of 480 multi-turn conversations (240 Crescendo-style jailbreak, 240 benign from WildChat). An additional 65 FITD and 58 Deceptive Delight conversations are generated for cross-attack evaluation. Test set contains 72 conversations. Hard negatives are sourced from 120 WildChat conversations exhibiting high topic drift (filtered for contamination, final rate 3\%).

\section{Experiments}

\subsection{Experiment 1: Lexical Confound Audit}

Table~\ref{tab:tfidf_audit} shows TF-IDF separation at each normalization stage.

\begin{table}[h]
\centering
\caption{TF-IDF Separation After Each Normalization Round}
\label{tab:tfidf_audit}
\begin{tabular}{lcc}
\toprule
\textbf{Stage} & \textbf{TF-IDF F1} & \textbf{Observation} \\
\midrule
Raw data & 1.00 & Perfect lexical separation \\
Round 1 (paraphrase) & 0.82 & Template words removed \\
Round 2 (targeted) & 0.73 & Passes gate ($<0.75$) \\
\bottomrule
\end{tabular}
\end{table}

The raw data exhibits perfect lexical separation, meaning any reported detection accuracy on unnormalized data is uninterpretable. Two rounds of normalization are required to reduce TF-IDF below our gate threshold. Even after normalization, 47 features retain discriminative power (Table~\ref{tab:tfidf_clean}), indicating residual topic-level lexical signal that cannot be removed without destroying content validity.

\begin{table}[h]
\centering
\caption{TF-IDF + LR Performance on Clean (Normalized) Data}
\label{tab:tfidf_clean}
\begin{tabular}{lc}
\toprule
\textbf{Metric} & \textbf{Value} \\
\midrule
F1 & 0.70 \\
Precision & 0.72 \\
Recall & 0.68 \\
Discriminative features & 47 \\
\bottomrule
\end{tabular}
\end{table}

\subsection{Experiment 2: Treatment vs. Baseline on Clean Data}

Table~\ref{tab:core} compares the persuasion-grounded treatment against the vanilla semantic baseline across conversation prefixes of length $K$.

\begin{table}[h]
\centering
\caption{F1 Comparison on Clean Data (N=72 test). None of the differences are statistically significant.}
\label{tab:core}
\begin{tabular}{lcccl}
\toprule
\textbf{Prefix} & \textbf{Treatment} & \textbf{Baseline} & \textbf{$\Delta$} & \textbf{$p$-value} \\
\midrule
$K=1$ & 0.68 & 0.66 & +2pp & 0.41 \\
$K=2$ & 0.72 & 0.70 & +2pp & 0.35 \\
$K=3$ & 0.78 & 0.76 & +2pp & 0.32 \\
$K=5$ & 0.83 & 0.82 & +1pp & 0.58 \\
Full & 0.86 & 0.85 & +1pp & 0.62 \\
\bottomrule
\end{tabular}
\end{table}

The persuasion-grounded encoder provides at most 2 pp improvement, none statistically significant. The claimed early detection advantage (at $K=2$-$3$) does not materialize on clean data.

\subsection{Experiment 3: Encoder Quality Degradation}

Table~\ref{tab:encoder} shows that the DeBERTa persuasion classifier itself degrades substantially on normalized data.

\begin{table}[h]
\centering
\caption{DeBERTa Encoder Accuracy: Original vs. Clean Data}
\label{tab:encoder}
\begin{tabular}{lcc}
\toprule
\textbf{Task} & \textbf{Original} & \textbf{Clean} \\
\midrule
Persuasion strategy (8-class) & 0.796 & 0.63 \\
Intent classification & 0.951 & 0.88 \\
\bottomrule
\end{tabular}
\end{table}

Persuasion classification accuracy drops from 0.796 to 0.63 on clean data, while intent classification remains relatively stable (0.951 to 0.88). The 8-class taxonomy becomes unreliable after register normalization, as many strategy boundaries (e.g., ``authority appeal'' vs. ``logical appeal'') depend on surface linguistic cues that normalization removes.

\subsection{Experiment 4: Scrambled Label Ablation}

To test whether persuasion labels carry genuine signal, we retrain the encoder with randomly permuted strategy labels.

\begin{table}[h]
\centering
\caption{Scrambled Label Ablation (Clean Data)}
\label{tab:scrambled}
\begin{tabular}{lccc}
\toprule
\textbf{Prefix} & \textbf{Proper Labels} & \textbf{Scrambled} & \textbf{$\Delta$} \\
\midrule
$K=2$ & 0.72 & 0.70 & $-$2pp \\
$K=3$ & 0.78 & 0.76 & $-$2pp \\
Full & 0.86 & 0.84 & $-$2pp \\
\bottomrule
\end{tabular}
\end{table}

Scrambled labels cause only a 2 pp drop, well below any meaningful threshold. Correct persuasion labels are not a significant signal source in the current data regime. The encoder functions primarily as a regularized feature extractor regardless of label quality.

\subsection{Experiment 5: Adversarial Perturbation}

We apply aggressive paraphrasing to jailbreak conversations to test robustness.

\begin{table}[h]
\centering
\caption{Adversarial Perturbation Results (Full Conversation)}
\label{tab:adversarial}
\begin{tabular}{lccr}
\toprule
\textbf{Model} & \textbf{Before} & \textbf{After} & \textbf{$\Delta$} \\
\midrule
Treatment & 0.86 & 0.71 & $-$15pp \\
Baseline & 0.85 & 0.72 & $-$13pp \\
TF-IDF+LR & 0.73 & 0.58 & $-$15pp \\
\bottomrule
\end{tabular}
\end{table}

The delta-of-deltas is $-2$ pp (treatment degrades \emph{more} than baseline), directionally opposite to the pragmatic evasion asymmetry hypothesis. Both neural models lose approximately 14 pp under perturbation, indicating that neither has learned features more robust than surface-level linguistic patterns.

\subsection{Experiment 6: Cross-Attack Generalization}

Models trained on Crescendo are evaluated zero-shot on FITD (65 conversations, avg 6.2 turns) and Deceptive Delight (58 conversations, avg 4.8 turns).

\begin{table}[h]
\centering
\caption{Cross-Attack Zero-Shot Evaluation}
\label{tab:cross}
\begin{tabular}{lcccc}
\toprule
\textbf{Attack Type} & \textbf{Treatment} & \textbf{Baseline} & \textbf{$\Delta$} & \textbf{Sig.} \\
\midrule
FITD & 0.58 & 0.56 & +2pp & n.s. \\
Deceptive Delight & 0.54 & 0.52 & +2pp & n.s. \\
Average & 0.56 & 0.54 & +2pp & n.s. \\
\bottomrule
\end{tabular}
\end{table}

Both systems collapse to near-chance performance on unseen attack types (from 0.85-0.86 in-domain to 0.54-0.58 cross-attack). Persuasion grounding provides no generalization benefit, contradicting the hypothesis that it should capture attack-agnostic pragmatic escalation patterns.

\subsection{Experiment 7: Hard Negative Evaluation}

We evaluate false positive rates on 120 WildChat conversations exhibiting high topic drift but no adversarial intent.

\begin{table}[h]
\centering
\caption{False Positive Rate on Hard Negatives}
\label{tab:fpr}
\begin{tabular}{lcc}
\toprule
\textbf{Condition} & \textbf{Treatment FPR} & \textbf{Baseline FPR} \\
\midrule
Hard negatives (high drift) & 0.22 & 0.20 \\
Easy negatives & 0.08 & 0.09 \\
\bottomrule
\end{tabular}
\end{table}

Both models misclassify 20-22\% of high-drift benign conversations as jailbreaks. Persuasion grounding does not help distinguish adversarial escalation from benign conversational drift. These false positive rates are unacceptable for deployment.

\subsection{Experiment 8: Strategy Transition Analysis}

We analyze persuasion strategy transition patterns in jailbreak vs. benign conversations.

\begin{table}[h]
\centering
\caption{Trajectory Statistics: Jailbreak vs. Benign}
\label{tab:trajectory}
\begin{tabular}{lccc}
\toprule
\textbf{Metric} & \textbf{Jailbreak} & \textbf{Benign} & \textbf{$\Delta$} \\
\midrule
Monotonicity score & 0.52 & 0.45 & 0.07 \\
Transition entropy & 2.1 & 2.3 & 0.2 \\
Dominant path freq. & 0.18 & 0.14 & 0.04 \\
\bottomrule
\end{tabular}
\end{table}

The hypothesized monotonic escalation trajectory is not observed. Jailbreak conversations show only marginally higher monotonicity (0.52 vs. 0.45) and marginally lower transition entropy (2.1 vs. 2.3). These differences are too small to serve as a detection signal, likely because the persuasion classifier's low accuracy (0.63) on normalized data produces noisy strategy assignments.

\subsection{Experiment 9: Cross-Checkpoint Variance}

\begin{table}[h]
\centering
\caption{Cross-Checkpoint Evaluation}
\label{tab:checkpoint}
\begin{tabular}{lcc}
\toprule
\textbf{Checkpoint} & \textbf{Persuasion Acc} & \textbf{Detection F1} \\
\midrule
Epoch 2 & 0.59 & 0.84 \\
Epoch 3 & 0.63 & 0.86 \\
Epoch final & 0.62 & 0.85 \\
\midrule
Variance & -- & 2pp \\
\bottomrule
\end{tabular}
\end{table}

Detection performance varies by only 2 pp across checkpoints despite persuasion accuracy ranging from 0.59 to 0.63. The weak correlation between encoder classification quality and downstream detection further confirms that persuasion labels are not the active ingredient.

\subsection{Experiment 10: DeepContext Proxy Comparison}

\begin{table}[h]
\centering
\caption{DeepContext Proxy (BERT-base + GRU) vs. Treatment}
\label{tab:deepcontext}
\begin{tabular}{lccc}
\toprule
\textbf{Prefix} & \textbf{Proxy} & \textbf{Treatment} & \textbf{$\Delta$} \\
\midrule
$K=2$ & 0.69 & 0.72 & +3pp \\
$K=3$ & 0.75 & 0.78 & +3pp \\
Full & 0.83 & 0.86 & +3pp \\
\bottomrule
\end{tabular}
\end{table}

The treatment exceeds the DeepContext proxy by a consistent 3 pp, matching the expected capacity gap between DeBERTa-v3-base and BERT-base on standard NLU benchmarks (typically 2-4 pp). This suggests the improvement is attributable to encoder capacity rather than persuasion grounding.

\section{Discussion}

\paragraph{Lexical confounds are structural, not incidental.}
The TF-IDF=1.0 finding is not a data quality bug specific to one dataset. Multi-turn attack data is necessarily generated from prompt templates (researchers cannot ethically collect real attacks at scale), while benign data comes from naturalistic user interactions. This asymmetry in data provenance creates register differences that any lexical classifier can exploit. Even after two rounds of aggressive normalization, 47 lexical features retain discriminative power (TF-IDF F1=0.70 on ``clean'' data). Complete deconfounding without destroying content validity may be impossible, raising questions about what any detector trained on such data actually learns.

\paragraph{The evasion hierarchy collapses.}
Our experiments reveal the detection performance hierarchy: lexical features (F1=0.70) $<$ semantic embeddings (F1=0.85) $\approx$ pragmatic features (F1=0.86). The theoretically motivated distinction between semantic and pragmatic levels does not manifest empirically. Both neural approaches degrade identically under perturbation ($-$13 to $-$15 pp) and generalize identically to unseen attacks (F1=0.54-0.58). Under current taxonomy granularity and data scale, the pragmatic level offers no advantage over the semantic level.

\paragraph{Implications for the field.}
These results suggest that published multi-turn detection accuracies should be interpreted with caution unless accompanied by a lexical confound audit. We recommend that future work report TF-IDF baseline performance alongside neural detector results as standard practice. The register-normalization protocol we propose provides a practical diagnostic: if TF-IDF separation exceeds 0.75, reported improvements may reflect shortcut learning rather than genuine capability.

\paragraph{Limitations.}
Our test set (N=72) limits statistical power. The register-normalization protocol itself may introduce artifacts. The 8-class persuasion taxonomy may be too fine-grained for the 480-sample data regime; coarser taxonomies (3-4 classes) or larger-scale annotation could yield different conclusions.

\section{Conclusion}

We present a methodology for auditing lexical confounds in multi-turn jailbreak detection benchmarks. Our register-normalization protocol reveals that a representative persuasion-grounded detector, which appears highly effective on raw data, provides no significant advantage over vanilla semantic baselines once lexical shortcuts are controlled. The persuasion labels themselves carry negligible signal in the current data regime (scrambled labels cause only 2 pp degradation). Cross-attack generalization and adversarial robustness are equivalently poor for both approaches. These findings argue that the multi-turn jailbreak detection community needs standardized confound auditing before claiming detection capability. We release our protocol to facilitate such audits.

\bibliographystyle{plain}
\begin{thebibliography}{10}

\bibitem{bullwinkel2025}
Bullwinkel, B., et al. (2025).
Multi-turn jailbreak attacks exploit per-turn representation safety.
\textit{arXiv preprint arXiv:2507.02956}.

\bibitem{zeng2024pap}
Zeng, Y., et al. (2024).
How Johnny Can Persuade LLMs to Jailbreak Them: Rethinking Persuasion to Challenge AI Safety by Humanizing LLMs.
\textit{Proceedings of ACL 2024}.

\bibitem{hagendorff2026}
Hagendorff, T. (2026).
Persuasion strategies in LLM jailbreak attacks.
\textit{Nature Communications}.

\bibitem{austin1962}
Austin, J. L. (1962).
\textit{How to Do Things with Words}.
Oxford University Press.

\bibitem{searle1969}
Searle, J. R. (1969).
\textit{Speech Acts: An Essay in the Philosophy of Language}.
Cambridge University Press.

\bibitem{deepcontext2026}
DeepContext (2026).
Contextual drift detection for multi-turn safety.
\textit{Preprint}.

\bibitem{gonen2019}
Gonen, H. and Goldberg, Y. (2019).
Lipstick on a Pig: Debiasing Methods Cover up Systematic Gender Biases in Word Embeddings But do not Remove Them.
\textit{Proceedings of NAACL-HLT 2019}.

\bibitem{persuade2025}
How Do LLMs Persuade? (2025).
Persuasion dynamics in turn-level representations.
\textit{arXiv preprint arXiv:2508.05625}.

\end{thebibliography}

\end{document}
